\newcommand{\prom}[1]{\left\langle #1 \right\rangle}
\newcommand{\norm}[1]{\left\|#1 }
\renewcommand{\AA}{\mathcal{A}}
\newcommand{\BB}{\mathcal{B}}
\newcommand{\CC}{\mathcal{C}}
\newcommand{\DD}{\mathcal{D}}
\newcommand{\II}{\mathcal{I}}
\newcommand{\kx}{K\langle X\rangle}

\newcommand{\name}[1]{ #1_\CC}

\newcommand{\setting}{\mathfrak S}
\newcommand{\events}{\mathfrak E}
\newcommand{\maxevents}{\mathfrak m}
\newcommand{\vsp}[1]{V_{#1}}

\newcommand{\pmodel}{\mathcal P}
\newcommand{\pnons}{\mathcal N}
\newcommand{\plocal}{\mathcal L}
\newcommand{\pquant}{\mathcal Q}

\newcommand{\disjoint}{D(S)}
\newcommand{\opa}{\mathfrak A}
\def \xxx {\langle \vec x \rangle}
\def \sxx {\xx [\rho]}
\def \txx {\xx [\tau]}
\renewcommand{\ss}[1]{M(#1)}

\def \P {\mathfrak P}
\def \O {\mathfrak O}

\newcommand{\zero}{\mathbf{0}}
\newcommand{\one}{\mathbf{1}}

\newcommand{\parts}[1]{\mathcal P(#1)}
\renewcommand{\norm}[1]{\left \| #1 \right\|}
\newcommand{\innerprod}[2]{\langle #1, #2 \rangle}

\newcommand{\st}{\operatorname{s.t.}}

% matrices
\newcommand{\matrices}{\operatorname{Mat}}
\newcommand{\rank}{\operatorname{rank}}
% unit ball
\newcommand{\ball}[1]{\operatorname{Ball}(#1)}
% extremal
\newcommand{\ext}{\operatorname{ext}}
% support
\newcommand{\supp}[1]{\operatorname{supp}(#1)}
% convex hull
\newcommand{\conv}{{\operatorname{conv}}}

% ame state
\newcommand{\ame}[2]{\operatorname{AME}(#1, #2)}

% Number fields
\newcommand{\N}{\mathds N}
\newcommand{\Q}{\mathds Q}
\newcommand{\R}{\mathds R}
\newcommand{\C}{\mathds C}
% Grobner basis generated from v
\newcommand{\GB}[1]{\text{GB}\langle #1 \rangle}
% field
\newcommand{\kk}{\mathds K}
% complex conjugate
\newcommand{\conj}[1]{#1^\inv}
% adjoint
\newcommand{\adj}{*}
% Letters
\newcommand{\xx}{X}
\newcommand{\pp}{P}
% Presentation
\newcommand{\pres}[2]{\langle #1 ; #2 \rangle}
% Monoid
\newcommand{\monoid}{\mathds M}
\newcommand{\ncmonoid}{\mathcal W}
\newcommand{\ncalgebra}{\mathcal P}
\newcommand{\pcalgebra}{\kk \ncmonoid}
\newcommand{\gpalgebra}{\mathcal{GP}}
\newcommand{\cmonoid}{[ \xx ]}
\newcommand{\pcmonoid}{\pres{\xx}{C}}
\newcommand{\statemonoid}{\mathds M_S(\xx, \rho)}
\newcommand{\tracemonoid}{\mathds M_T(\xx, \tau)}
\newcommand{\formalmonoid}{\llangle \mathds M \rrangle}
\newcommand{\lpmonoid}{[\xx|\pp]}
% Involution
\newcommand{\inv}{*}
% Polynomial rings
\newcommand{\ncring}{k \ncmonoid }
\newcommand{\cring}{k \cmonoid}
\newcommand{\pcring}{k \pcmonoid }
\newcommand{\statering}{\kk \langle \statemonoid \rangle}
\newcommand{\tracering}{\kk \langle \tracemonoid \rangle}
\newcommand{\lpring}{\kk \lpmonoid }
\newcommand{\polyalg}{R}
\newcommand{\monoidring}{\kk \langle \monoid \rangle}

% Equivalence classes
\newcommand{\canonical}[1]{ \left\ulcorner #1 \right\urcorner }
\newcommand{\eqclass}[1]{ [ #1]}
\newcommand{\replacement}[1]{\mapsto_{#1}}
\newcommand{\ideal}{\mathcal I}
\newcommand{\idealgen}[1]{\langle #1 \rangle}
\newcommand{\monoidgen}[1]{\langle #1 \rangle}
% Leading terms
\newcommand{\lc}[1]{\operatorname{lc}(#1)}
\newcommand{\lm}[1]{\operatorname{lm}(#1)}
\newcommand{\lt}[1]{\operatorname{lt}(#1)}
% Letterplace
\newcommand{\shift}[1]{\operatorname{sh}(#1)}
%  maximal-clique
\newcommand{\clique}{\mathcal C}
% Jordan algebra
% linear space
\newcommand{\LL}{\mathcal L}
% linear span
\newcommand{\lspn}{\operatorname{span}}

% expectation
\newcommand{\expectation}[1]{\langle #1 \rangle}

% clifford algebra
\newcommand{\clifford}{\mathfrak C}

% moments
\newcommand{\moments}{\mathcal M}

% highlight comments
\newcommand{\moi}[1]{\textcolor{yellow!70!red!80!black}{#1}}

\newcommand{\AM}[1]{\textcolor{blue!80!black}{#1}}
\newcommand{\aM}[1]{\textcolor{purple!80!black}{#1}}

\newcommand{\SP}[1]{\textcolor{red!80!black}{#1}}

% multiplication
\newcommand{\mult}{ }
% involution
\newcommand{\invol}{*}

% Color boxes
\colorlet{summarycolor}{yellow!70!red!90!blue!80!white}
\colorlet{algorithmcolor}{blue!50!white}
\colorlet{softwarecolor}{red!50!white}

\newtcbtheorem[auto counter,number within=section]{summary}%
  {Highlight}{fonttitle=\bfseries\upshape, fontupper=\slshape,
     arc=0mm, colback=blue!5!white,colframe=summarycolor}{theorem}
     
\newcommand{\summarybox}[2]{
\begin{figure}[h!]
\begin{tcolorbox}[colback=summarycolor!70!white,colframe=summarycolor,title=#1, fonttitle=\bfseries]
  \raggedright{#2}
\end{tcolorbox}
\end{figure}
}

\newcommand{\algorithmbox}[2]{
\begin{figure}[h!]
\begin{tcolorbox}[colback=white,colframe=algorithmcolor,title=#1, fonttitle=\bfseries]
  #2
\end{tcolorbox}
\end{figure}
}

\newcommand{\softwarebox}[2]{
\begin{figure}[h!]
\begin{tcolorbox}[colback=softwarecolor!70!white,colframe=softwarecolor,title=#1, fonttitle=\bfseries]
  #2
\end{tcolorbox}
\end{figure}
}

\newcommand\blueout{\bgroup\markoverwith
{\textcolor{blue}{\rule[.5ex]{2pt}{0.4pt}}}\ULon}
\newcommand{\mathsout}[1]% will draw line through middle of #1
{\bgroup\mathchoice
  {\sbox0{$\displaystyle{#1}$}%
    \usebox0\hspace{-\wd0}%
    \rule[0.5\ht0-0.5\dp0-.5pt]{\wd0}{1pt}}%
  {\sbox0{$\textstyle{#1}$}%
    \usebox0\hspace{-\wd0}%
    \rule[0.5\ht0-0.5\dp0-.5pt]{\wd0}{1pt}}%
  {\sbox0{$\scriptstyle{#1}$}%
    \usebox0\hspace{-\wd0}%
    \rule[0.5\ht0-0.5\dp0-.5pt]{\wd0}{1pt}}%
  {\sbox0{$\scriptscriptstyle{#1}$}%
    \usebox0\hspace{-\wd0}%
    \rule[0.5\ht0-0.5\dp0-.5pt]{\wd0}{1pt}}%
\egroup}
\let \vec \textbf
\newcommand{\id}{\mathbf{1}}

\newcommand{\ketbra}[2]{\mathinner{|#1\rangle \langle#2|}}
\newcommand{\dyad}[1]{| #1\rangle \langle #1|}
\newcommand\hcancel[2][black]{\setbox0=\hbox{$#2$}%
\rlap{\raisebox{.45\ht0}{\textcolor{#1}{\rule{\wd0}{1pt}}}}#2} 
\newcommand{\nn}{\nonumber}
\newcommand{\tr}{\operatorname{tr}}

\newcommand{\adjpair}{\pi}


% Tracial optimization
\newcommand{\inonto}{\lhook\joinrel\twoheadrightarrow}
\newcommand{\onto}{\twoheadrightarrow}
\newcommand{\into}{\hookrightarrow}
\newcommand{\headparagraph}[1]{\noindent \textbf{#1} --}
\newcommand{\FF}{\mathds{F}}
\newcommand{\Hclass}[1]{\llbracket #1 \rrbracket}
\newcommand{\Gcan}[1]{\lceil #1 \rceil}
\newcommand\cyclic[1]{\mathring{#1}}
\renewcommand{\AA}{\mathcal{A}}
\newcommand{\EE}{\mathcal{E}}
\newcommand{\HH}{\mathcal{H}}
\newcommand{\PP}{\mathcal{P}}
\newcommand{\JJ}{\mathcal{J}}

\newcommand{\csupp}{\operatorname{csupp}}

%vectores en negrita
\let\vec\mathbf
%epsilon
\let\epsilon\varepsilon


\newcommand{\graph}{\mathcal G}
%\newcommand{\canonical}[1]{[#1]}
\newcommand{\complem}[1]{\bar #1}
\newcommand{\pcmonoidgraph}{\ncmonoid_\graph}
\newcommand{\pcalgebragraph}{\ncalgebra_\graph}

\newcommand{\cliqueword}[1]{#1^\clique}
\newcommand{\graphword}[1]{#1^G}

\newcommand{\graphcanonical}[1]{[#1]}

\newcommand{\sos}{\operatorname{SOS}_{R, Q}}